%-------------------------------------------------------------------------------
% Bartha István Zénó – Önéletrajz (magyar változat)
% Europass szakaszszerkezet, "Jake's Resume" vizuális stílus.
% Fordítás: pdflatex / tectonic BarthaZeno_CV_HU.tex  (kell hozzá: zeno_cv_small.jpg)
%-------------------------------------------------------------------------------
\documentclass[a4paper,11pt]{article}

\usepackage[T1]{fontenc}
\usepackage[utf8]{inputenc}
\usepackage{lmodern}
\usepackage[empty]{fullpage}
\usepackage{titlesec}
\usepackage[dvipsnames]{xcolor}
\usepackage{enumitem}
\usepackage{hyperref}
\usepackage{fancyhdr}
\usepackage{tabularx}
\usepackage{array}
\usepackage{fontawesome5}
\usepackage{tikz}
\usepackage{graphicx}

%---------------------------- Colours & links ----------------------------------
\definecolor{accent}{RGB}{28,63,110}   % deep navy
\definecolor{muted}{RGB}{90,90,90}
\hypersetup{colorlinks=true, urlcolor=accent, linkcolor=accent,
  pdftitle={Bartha István Zénó – Önéletrajz}, pdfauthor={Bartha István Zénó}}

%---------------------------- Page layout --------------------------------------
\pagestyle{fancy}
\fancyhf{}
\renewcommand{\headrulewidth}{0pt}
\renewcommand{\footrulewidth}{0pt}

\addtolength{\oddsidemargin}{-0.55in}
\addtolength{\evensidemargin}{-0.55in}
\addtolength{\textwidth}{1.1in}
\addtolength{\topmargin}{-0.65in}
\addtolength{\textheight}{1.3in}

\urlstyle{same}
\raggedbottom
\raggedright
\setlength{\tabcolsep}{0in}
\pdfgentounicode=1

%---------------------------- Sections -----------------------------------------
\titleformat{\section}
  {\vspace{-4pt}\color{accent}\scshape\raggedright\large}
  {}{0em}{}
  [\color{accent}\titlerule\vspace{-5pt}]

%---------------------------- Custom commands ----------------------------------
\newcommand{\resumeItem}[1]{\item\small{#1\vspace{-2pt}}}

% {Title}{Right top}{Subtitle}{Right bottom}
\newcommand{\resumeSubheading}[4]{
  \vspace{-2pt}\item
    \begin{tabular*}{0.97\textwidth}[t]{l@{\extracolsep{\fill}}r}
      \textbf{#1} & #2 \\
      \textit{\small #3} & \textit{\small #4} \\
    \end{tabular*}\vspace{-7pt}
}

% {Title | stack}{Right label}
\newcommand{\resumeProjectHeading}[2]{
  \item
    \begin{tabular*}{0.97\textwidth}{l@{\extracolsep{\fill}}r}
      \small #1 & \textit{\small\color{muted} #2} \\
    \end{tabular*}\vspace{-7pt}
}

% Clickable GitHub icon next to a project title: \ghlink{repo-name}
\newcommand{\ghuser}{juanmiguelaasg}
\newcommand{\ghlink}[1]{\,\href{https://github.com/\ghuser/#1}{\faGithub}}

\renewcommand\labelitemii{$\vcenter{\hbox{\tiny$\bullet$}}$}
\newcommand{\resumeSubHeadingListStart}{\begin{itemize}[leftmargin=0.15in, label={}]}
\newcommand{\resumeSubHeadingListEnd}{\end{itemize}}
\newcommand{\resumeItemListStart}{\begin{itemize}[itemsep=1pt]}
\newcommand{\resumeItemListEnd}{\end{itemize}\vspace{-5pt}}

%===============================================================================
\begin{document}

%--------------------------------- HEADER --------------------------------------
\noindent
\begin{minipage}[c]{3.1cm}
  \begin{tikzpicture}
    \clip (0,0) circle (1.45cm);
    % Portrait photo (3:4) shifted down so the face sits in the centre of the circle.
    \node at (0,-0.45cm) {\includegraphics[width=3.2cm]{zeno_cv_small.jpg}};
  \end{tikzpicture}
\end{minipage}%
\hspace{0.5cm}%
\begin{minipage}[c]{\dimexpr\textwidth-3.6cm\relax}
  {\Huge\scshape\color{accent} Bartha István Zénó} \\[3pt]
  {\large Üzemmérnök-informatikus hallgató \,$\cdot$\, Workflow automatizáció \& AI} \\[5pt]
  \small
  \faMobile*~+36 20 263 7679 \quad
  \href{mailto:zenoka14@gmail.com}{\faEnvelope~zenoka14@gmail.com} \quad
  \faMapMarker*~Budapest \\[2pt]
  \href{https://www.linkedin.com/in/z\%C3\%A9n\%C3\%B3-bartha-a9814b354/}{\faLinkedin~LinkedIn} \quad
  \href{https://github.com/juanmiguelaasg}{\faGithub~github.com/juanmiguelaasg} \quad
  \href{https://juanmiguelaasg.github.io/portfolio/}{\faGlobe~juanmiguelaasg.github.io/portfolio}
\end{minipage}

\vspace{4pt}

%--------------------------------- ABOUT ME ------------------------------------
\section{Bemutatkozás}
\small
Üzemmérnök-informatikus hallgató vagyok az Óbudai Egyetemen. Őszintén érdekel, hogyan lehet az IT-rendszereket
megbízhatóan működtetni -- és hogyan tehetik ezt a munkát egyszerűbbé az automatizáció és az AI. Szeretem a problémákat
az alapjaiktól megérteni, gyorsan tanulok, és szívesen alakítom az ismétlődő feladatokat átlátható, jól felépített
munkafolyamatokká. Olyan lehetőséget keresek egy Operation teamben, ahol hozzá tudok tenni a csapat munkájához,
tanulhatok tapasztalt kollégáktól, és valódi vállalati környezetben fejlődhetek.
\normalsize

%-------------------------- EDUCATION AND TRAINING -----------------------------
\section{Tanulmányok}
  \resumeSubHeadingListStart
    \resumeSubheading
      {Óbudai Egyetem}{Budapest}
      {BProf Üzemmérnök-informatikus \,--\, \textbf{Tanulmányi átlag: 4,2/5,0}}{2024. szept. -- jelenleg}
      \resumeItemListStart
        \resumeItem{\textbf{Tanulmányi ösztöndíj} a 2. és a 3. félévben}
        \resumeItem{\textbf{Releváns tárgyak:} Számítógép-hálózatok, Adatbázisok, Beágyazott és érzékelőrendszerek,
          Rendszerelmélet, Szoftverfejlesztés alapjai, Problémamegoldás programozással, Elektronika, Vállalkozás és
          projektmenedzsment, Diszkrét matematika, Analízis}
      \resumeItemListEnd
    \resumeSubheading
      {Magyar Agrár- és Élettudományi Egyetem (MATE)}{Budapest}
      {Élelmiszermérnök BSc tanulmányok -- nem fejeztem be; szakmát váltottam az IT irányába}{2021 -- 2024}
    \resumeSubheading
      {Szent Margit Gimnázium}{Budapest}
      {Érettségi}{2017 -- 2021}
      \resumeItemListStart
        \resumeItem{\textbf{Emelt szint:} Angol \quad $\cdot$ \quad \textbf{Középszint:} Matematika, Informatika}
      \resumeItemListEnd
    \resumeSubheading
      {Újbudai Ádám Jenő Általános Iskola}{Budapest}
      {Általános iskola}{2009 -- 2017}
  \resumeSubHeadingListEnd

%---------------------------------- PROJECTS -----------------------------------
\section{Projektek}
  \resumeSubHeadingListStart

    \resumeProjectHeading
      {\textbf{System Process Monitor}\ghlink{SystemProcessMonitor} $|$ \emph{Python, SQL, idősorok, anomáliadetektálás, webhook, Docker}}{Saját projekt}
      \resumeItemListStart
        \resumeItem{Könnyűsúlyú, lokálisan futó \textbf{Python monitorozó ágens}, amely rendszer- és folyamatszintű metrikákat
          gyűjt (CPU, memória, lemez, hálózat, folyamatok és szálak), és \textbf{idősorként SQL-ben (SQLite)} tárolja őket.}
        \resumeItem{\textbf{Anomáliadetektálás} gördülő baseline-nal (EWMA + z-score) fix küszöbök helyett, így a gép megszokott
          terhelése nem vált ki riasztást -- csak a hirtelen eltérés.}
        \resumeItem{A duplikációszűrt riasztásokat \textbf{incidensként továbbítja webhookon / REST-en} (megnyitás és lezárás,
          stabil dedup kulcs, újrapróbálás), így egy ticketkezelő rendszer fel tudja dolgozni őket.}
        \resumeItem{Asztali felület \textbf{élő dashboarddal}, SQL-riportokkal és szabályszerkesztővel; az ágens ablak nélkül és
          \textbf{Dockerben} is futtatható.}
      \resumeItemListEnd

    \resumeProjectHeading
      {\textbf{Automated API Health-Checker}\ghlink{api-health-checker} $|$ \emph{Python, REST, YAML, Azure DevOps}}{Saját projekt}
      \resumeItemListStart
        \resumeItem{Ütemezetten futó ellenőrző, amely a YAML-ban leírt \textbf{REST végpontokat} teszteli -- státuszkód, válaszidő
          és \textbf{JSON-séma} --; egy \textbf{Azure DevOps pipeline} lépéseként is fut.}
        \resumeItem{Újrapróbálás és megerősítési logika szűri a téves riasztásokat; végpontonként méri a \textbf{rendelkezésre
          állást az SLA-célhoz}, és egyszerű státuszoldalt készít.}
      \resumeItemListEnd

    \resumeProjectHeading
      {\textbf{AI Log Analyzer}\ghlink{ai-log-analyzer} $|$ \emph{Python, LLM, RAG}}{Saját projekt}
      \resumeItemListStart
        \resumeItem{Több forrásból érkező logokat dolgoz fel és egységesít, majd \textbf{klaszterezi az ismétlődő hibákat}, így
          több száz sorból néhány jól elkülönülő probléma lesz.}
        \resumeItem{Egy \textbf{LLM a belső dokumentációban keresve} összefoglalja az egyes problémákat, valószínű kiváltó okot
          javasol és a releváns runbookra mutat; a \textbf{konfidencia-küszöb} alatti eredmények emberi felülvizsgálatra kerülnek.}
      \resumeItemListEnd

    \resumeProjectHeading
      {\textbf{API Integration Hub}\ghlink{api-integration-hub} $|$ \emph{Python, REST, SQLite, openpyxl, CustomTkinter}}{Saját projekt}
      \resumeItemListStart
        \resumeItem{Az API Health-Checkerem továbbfejlesztése desktop alkalmazásként: \textbf{CRM, ERP és Helpdesk rendszereket
          köt össze REST-en} (API key / Bearer / Basic hitelesítés), YAML-ban leírt \textbf{mezőleképezéssel}; csak az új vagy
          megváltozott rekordokat küldi, így nincs duplikátum, a \textbf{health gate} pedig leállítja a szinkront, ha egy
          rendszer nem elérhető.}
        \resumeItem{Beépített \textbf{Scrum} (backlog, sprintek, tábla, burndown, velocity) -- a magas prioritású ticketek
          automatikusan backlog-storyk lesznek -- és egy kattintásos \textbf{Excel-riport}; 27 pytest teszt, GitHub Actions CI.}
      \resumeItemListEnd

    \resumeProjectHeading
      {\textbf{IT Kft. vállalkozási és projektterve}}{Egyetemi projekt}
      \resumeItemListStart
        \resumeItem{Egy kitalált IT-szolgáltató Kft. teljes tervének elkészítése: \textbf{szolgáltatási portfólió},
          \textbf{működési folyamatok} (support, szállítás, ügyfél-onboarding), szervezet, \textbf{költségvetés},
          mérföldkövekre bontott \textbf{ütemterv} és \textbf{kockázati nyilvántartás}.}
      \resumeItemListEnd

    \resumeProjectHeading
      {\textbf{Pszeudo-3D raycasting motor a terminálban}\ghlink{DOOMgame} $|$ \emph{C\#, .NET}}{Egyetemi projekt}
      \resumeItemListStart
        \resumeItem{Belső nézetes, Doom által ihletett \textbf{raycasting} renderer egyszerű terminálablakban, pályabetöltéssel,
          mozgással, \textbf{ütközésvizsgálattal} és fix időlépéses játékciklussal.}
      \resumeItemListEnd

    \resumeProjectHeading
      {\textbf{Okos kuka} $|$ \emph{C, Arduino (regiszterszinten)}}{Egyetemi projekt}
      \resumeItemListStart
        \resumeItem{A mikrokontrollert \textbf{közvetlenül a regiszterein keresztül} programoztam (időzítők, GPIO,
          \textbf{megszakítások}) magas szintű könyvtárak helyett; szenzorvezérelt fedélnyitás és telítettségmérés.}
      \resumeItemListEnd

    \resumeProjectHeading
      {\textbf{Robotporszívó vezérlés} $|$ \emph{ROS2}}{Egyetemi projekt}
      \resumeItemListStart
        \resumeItem{\textbf{ROS2} node-ok és topicok a szenzoradatok feldolgozására és a sebességparancsok kiadására,
          akadálykezeléssel és területlefedési logikával, szimulációban tesztelve.}
      \resumeItemListEnd

  \resumeSubHeadingListEnd

%------------------------------- DIGITAL SKILLS --------------------------------
\section{Digitális készségek}
 \begin{itemize}[leftmargin=0.15in, label={}]
    \small{\item{
     \textbf{Programozás \& szkriptelés:} Python (requests, pandas, psutil, openpyxl), C\#, C, SQL, Bash, PowerShell, HTML, CSS \\
     \textbf{Rendszerek \& hálózatok:} Linux alapok (CLI, jogosultságok, SSH, systemd), hálózati alapok (TCP/IP, DNS,
       DHCP, HTTP), Docker alapok, Microsoft Azure alapok \\
     \textbf{IT-szolgáltatásmenedzsment:} ITIL fogalmak -- incident-, problem-, change- és request-kezelés; SLA; priorizálás \\
     \textbf{Monitoring \& observability:} metrikagyűjtés, baseline-ok, anomáliadetektálás, riasztáskorreláció;
       Grafana, Prometheus alapok \\
     \textbf{Automatizáció \& AI:} workflow-orchesztráció, runbook automatizáció, Power Automate, OpenAI / Azure OpenAI API,
       prompt engineering, RAG, human-in-the-loop tervezés \\
     \textbf{API-k \& tesztelés:} REST API-k, Postman, webhookok, JSON, YAML \\
     \textbf{Verziókezelés \& CI/CD:} Git, GitHub, CI/CD alapok (GitHub Actions, Azure DevOps Pipelines), Azure DevOps Boards \\
     \textbf{Irodai eszközök:} Microsoft 365 (Excel, Word, PowerPoint, Teams) \\
     \textbf{Beágyazott rendszerek \& robotika:} Arduino (regiszterszintű C), ROS2
    }}
 \end{itemize}

%------------------------------ PERSONAL SKILLS --------------------------------
\section{Személyes készségek}
 \begin{itemize}[leftmargin=0.15in, label={}]
    \small{\item{
     \textbf{Problémamegoldás:} szisztematikus hibakeresés egészen hardverszintig (regiszterszintű mikrokontroller-programozás) \\
     \textbf{Csapatmunka:} két hackathonon időkorlátos csapatban készítettünk működő prototípust, a feladatokat Azure DevOps-ban,
       Agile / Scrum alapelvek szerint koordinálva \\
     \textbf{Kommunikáció:} világos írásbeli és szóbeli kommunikáció magyarul és angolul; kórházi és könyvtári önkéntesként
       tapasztalat az emberek segítésében \\
     \textbf{Szervezés:} strukturált tervezés és dokumentáció -- ütemtervek, költségvetések és kockázati nyilvántartások
       (IT Kft. vállalkozási terv)
    }}
 \end{itemize}

%------------------------------ LANGUAGE SKILLS --------------------------------
\section{Nyelvtudás}
 \begin{itemize}[leftmargin=0.15in, label={}]
    \small{\item{
     \textbf{Magyar:} anyanyelv \\
     \textbf{Angol:} B2 -- Euroexam B2 nyelvvizsga (2020), emelt szintű érettségi
    }}
 \end{itemize}

%-------------------------- HACKATHONS & COMPETITIONS --------------------------
\section{Hackathonok és versenyek}
  \resumeSubHeadingListStart
    \resumeSubheading
      {GDE MIT Hackathon}{2026}
      {Résztvevő -- csapatalapú innovációs hackathon; a csapatmunkát \textbf{Azure DevOps}-ban szerveztük}{}
    \resumeSubheading
      {Byborg AI Hackathon}{2024}
      {Résztvevő -- a Byborg Enterprises AI-fókuszú fejlesztői versenye}{}
  \resumeSubHeadingListEnd

%------------------------ PROFESSIONAL DEVELOPMENT GOALS -----------------------
\section{Szakmai fejlődési céljaim}
  \begin{itemize}[leftmargin=0.3in, itemsep=1pt]
    \resumeItem{\textbf{Valódi vállalati tapasztalat:} megtapasztalni, hogyan működik az üzemeltetés nagyvállalati
      léptékben -- folyamatok, ügyelet és eszkaláció, change management, csapatok közötti együttműködés és üzleti hatás.
      Ezeket az egyetem be tudja mutatni, de igazán csak egy valódi üzemeltetési csapatban lehet megtanulni.}
    \resumeItem{\textbf{Üzemeltetési szakirány} az Óbudai Egyetemen -- az infrastruktúra-üzemeltetés, a
      szolgáltatásmenedzsment és az observability mélyebb megismerése, a szakdolgozatomat pedig az
      \textbf{AI-támogatott incidenskezelésről} szeretném írni.}
    \resumeItem{\textbf{Szakmai minősítések:} ITIL 4 Foundation és Microsoft Azure Fundamentals, majd további
      platformminősítések, ahogy gyakorlatot szerzek.}
    \resumeItem{\textbf{Nyelvek:} \textbf{német vagy spanyol B2 nyelvvizsga} megszerzése, hogy nemzetközi, többnyelvű
      csapatokban is hatékonyan dolgozzak.}
    \resumeItem{\textbf{Hosszú távon:} \textbf{automatizációs / site reliability mérnökké} válni, aki öngyógyító
      üzemeltetést épít -- ahol a rutint az automatizáció viszi, az emberek pedig a valóban nehéz problémákra
      koncentrálhatnak.}
  \end{itemize}
  \vspace{-4pt}

%--------------------------- ADDITIONAL INFORMATION ----------------------------
\section{További információk}
 \begin{itemize}[leftmargin=0.15in, label={}]
    \small{\item{
     \textbf{Önkéntesség:} kórházi önkéntes, Budapest (2021); könyvtári önkéntes, Budapest (2020 -- 2021) \\
     \textbf{Hobbik:} e-sport, főzés, túrázás, biciklizés, olvasás, sorozatok \\
     \textbf{Jogosítvány:} B kategória
    }}
 \end{itemize}

\end{document}
